\Needspace{8\baselineskip}
\subsection{General: formalidad y cumplimiento de instrucciones}
{\linespread{1}\LFXcuerpotabla
\setlength{\parskip}{0pt}\renewcommand{\arraystretch}{1.10}\setlength{\tabcolsep}{4pt}
\begin{longtable}{F{0.09} F{0.29} F{0.40} F{0.22}}
\caption[General: formalidad y cumplimiento de instrucciones]{General: formalidad y cumplimiento de instrucciones. Fuente: revisión del Informe 2 y versión corregida de LafroX.}\\[2pt]
\LFXtopruleoriginal \cab{N.º} & \cab{Observación} & \cab{Respuesta} & \cab{Sección modificada} \\ \midrule
\endfirsthead
\multicolumn{4}{@{}l@{}}{\fontsize{9pt}{11pt}\selectfont\color{lafroxGris}Tabla \thetable\ (continuación)}\\[3pt]
\LFXtopruleoriginal \cab{N.º} & \cab{Observación} & \cab{Respuesta} & \cab{Sección modificada} \\ \midrule
\endhead
\midrule
\multicolumn{4}{@{}r@{}}{\fontsize{9pt}{11pt}\selectfont\color{lafroxGris}Continúa en el folio siguiente.}\\
\endfoot
\bottomrule
\endlastfoot
\mbox{GEN-01} & Volumen excesivo: SD4 con 194 páginas de cuerpo y 117 de anexos; el subdocumento debe resumir y el anexo detallar. & \textbf{Parcial.} El SD4 traslada al Anexo 4-D los listados de M2, M3, M4 y M7, al 4-N los criterios de interacción y al 4-R la custodia. Se reducirá la extensión del cuerpo del SD4. El cuerpo del SD7 pasó de 116 a 90 páginas y sus formularios y anexos, de 376 a 324. Lo que falta se completa para la Propuesta Final. & SD4, 4.1.5 y Anexos 4-D, 4-N, 4-R; SD7, T-14, T-15, T-18 y Anexos del SD7 \\
\LFXfiletefila
\mbox{GEN-02} & La tabla del Art. 46° llega en un archivo fuera de la nomenclatura de la sección 2 de la Aclaración N.° 2. & \textbf{Corregida.} La tabla se entrega como Anexo A del Análisis económico-financiero (Entregable 2), en archivo propio con nombre conforme a la nomenclatura de la oferta. & Este anexo \\
\LFXfiletefila
\mbox{GEN-03} & Declaración de uso de IA con nivel «Alto» en todo y un único revisor; revisión humana de forma. & \textbf{Corregida.} Las Declaraciones de SD1, SD3, SD4, SD5, SD6 y SD8 asignan revisores distintos por sección, cada uno con el contenido que verificó; el SD13 agrega una pauta de revisión técnica. & Declaraciones de IA de SD1, SD3, SD4, SD5, SD6 y SD8; Anexos SD13 \\
\LFXfiletefila
\mbox{GEN-04} & Foliación única y correlativa de toda la propuesta. & \textbf{Pendiente.} Se numerarán los folios de forma correlativa, desde la portada del SD1 hasta el último formulario. Se corrige para la Propuesta Final. & Ensamble de la Propuesta Final \\
\LFXfiletefila
\mbox{GEN-05} & Erratas dentro de las figuras del SD2 («enERP», «enpapel», «retornarr», «No Diferencias no conciliadas»). & \textbf{Corregida.} Figuras reexportadas desde su fuente editable con el texto corregido; «No» queda como rótulo del conector. & SD2, Figuras 2.1, 2.5 y 2.6 \\
\LFXfiletefila
\mbox{GEN-06} & Sitio web corporativo en GitHub (lafrox.github.io) y correo con dominio lafrox.cl. & \textbf{Pendiente.} El dominio del sitio corporativo se unifica con el del correo. Se corrige para la Propuesta Final. & SD1, introducción \\
\end{longtable}
}

\Needspace{8\baselineskip}
\subsection{Subdocumento 1: Presentación de la empresa (T-6)}
{\linespread{1}\LFXcuerpotabla
\setlength{\parskip}{0pt}\renewcommand{\arraystretch}{1.10}\setlength{\tabcolsep}{4pt}
\begin{longtable}{F{0.09} F{0.29} F{0.40} F{0.22}}
\caption[Subdocumento 1: Presentación de la empresa (T-6)]{Subdocumento 1: Presentación de la empresa (T-6). Fuente: revisión del Informe 2 y versión corregida de LafroX.}\\[2pt]
\LFXtopruleoriginal \cab{N.º} & \cab{Observación} & \cab{Respuesta} & \cab{Sección modificada} \\ \midrule
\endfirsthead
\multicolumn{4}{@{}l@{}}{\fontsize{9pt}{11pt}\selectfont\color{lafroxGris}Tabla \thetable\ (continuación)}\\[3pt]
\LFXtopruleoriginal \cab{N.º} & \cab{Observación} & \cab{Respuesta} & \cab{Sección modificada} \\ \midrule
\endhead
\midrule
\multicolumn{4}{@{}r@{}}{\fontsize{9pt}{11pt}\selectfont\color{lafroxGris}Continúa en el folio siguiente.}\\
\endfoot
\bottomrule
\endlastfoot
\mbox{SD1-01} & Sin antecedentes financieros: ventas, patrimonio ni capacidad de garantía. & \textbf{Corregida.} Se agregan patrimonio 2023–2025, ventas 2025, liquidez corriente de 1,8 y capacidad de emisión de garantías. La acreditación documental se entrega en el Sobre 1. & SD1, 1.1 \\
\LFXfiletefila
\mbox{SD1-02} & El SOC aparece como existente, contratado y por implementar. & \textbf{Corregida.} Una sola versión: NOC y función CISO existentes; SOC propio en Valparaíso por habilitar con personal adicional, con calendario en el T-15. & SD1, 1.1, 1.2 y 1.5 \\
\LFXfiletefila
\mbox{SD1-03} & Demanda de desarrollo poco clara (31 FTE frente a 40 simultáneos y 48 disponibles); debe cuadrar con el T-15. & \textbf{Corregida.} La planta cuenta con 48 desarrolladores y 15 especialistas de arquitectura e integración. En el mes 7, el T-15 programa 34 de las 36 plazas del grupo de Desarrollo: una de dirección y 33 de construcción. & SD1, 1.2 y 1.5; T-15, 5.7 \\
\LFXfiletefila
\mbox{SD1-04} & El organigrama omite Comercial y Administración y no ubica los dos centros de datos. & \textbf{Corregida.} Nueva Figura 1.1 con Comercial y Administración y Finanzas bajo la Gerencia; centro primario en Santiago y secundario en Valparaíso; se agrega el área de Implantación y Gestión del Cambio (9 especialistas, dentro de una dotación técnica de 133). & SD1, 1.2, Figura 1.1 y Tabla 1.3 \\
\LFXfiletefila
\mbox{SD1-05} & La sección 1.5 nomina al equipo del proyecto, reservado al Capítulo 12. & \textbf{Corregida.} La 1.5 describe siete ámbitos y diez roles sin nombres y remite el equipo nominado al T-8. & SD1, 1.5 \\
\LFXfiletefila
\mbox{SD1-06} & El plan de renovación de certificaciones se presenta dos veces. & \textbf{Corregida.} Se elimina el párrafo repetido y se conservan las fichas por certificación; se corrige el plazo atribuido a ISO/IEC 20000-1. & SD1, 1.4 \\
\LFXfiletefila
\mbox{SD1-07} & No acredita el conocimiento sanitario y tributario de la RT-15.02 del caso. & \textbf{Corregida.} Se declara el conocimiento del Reglamento Sanitario de los Alimentos y de los DTE, respaldado en registros de capacitación y en los proyectos del T-6. & SD1, 1.4 \\
\LFXfiletefila
\mbox{SD1-08} & Las cuatro alianzas vencen entre los meses 2 y 11; falta analizarlo. & \textbf{Corregida.} Análisis del primer año: renovación 60 días antes del vencimiento, separación entre membresía y soporte y contingencia. & SD1, 1.6 \\
\LFXfiletefila
\mbox{SD1-09} & Referencia farmacéutica con cámaras a −22 °C poco habitual. & \textbf{Corregida.} Se precisa que son doce cámaras de reactivos y materiales de diagnóstico congelados. & SD1, 1.4, Tabla 1.5; T-6, proyecto 2 \\
\end{longtable}
}

\Needspace{8\baselineskip}
\subsection{Subdocumento 2: Comprensión del problema}
{\linespread{1}\LFXcuerpotabla
\setlength{\parskip}{0pt}\renewcommand{\arraystretch}{1.10}\setlength{\tabcolsep}{4pt}
\begin{longtable}{F{0.09} F{0.29} F{0.40} F{0.22}}
\caption[Subdocumento 2: Comprensión del problema]{Subdocumento 2: Comprensión del problema. Fuente: revisión del Informe 2 y versión corregida de LafroX.}\\[2pt]
\LFXtopruleoriginal \cab{N.º} & \cab{Observación} & \cab{Respuesta} & \cab{Sección modificada} \\ \midrule
\endfirsthead
\multicolumn{4}{@{}l@{}}{\fontsize{9pt}{11pt}\selectfont\color{lafroxGris}Tabla \thetable\ (continuación)}\\[3pt]
\LFXtopruleoriginal \cab{N.º} & \cab{Observación} & \cab{Respuesta} & \cab{Sección modificada} \\ \midrule
\endhead
\midrule
\multicolumn{4}{@{}r@{}}{\fontsize{9pt}{11pt}\selectfont\color{lafroxGris}Continúa en el folio siguiente.}\\
\endfoot
\bottomrule
\endlastfoot
\mbox{SD2-01} & El arbitraje ya adopta decisiones de solución (promesa 24/48 h con corte a las 14:00, retención por SSCC). & \textbf{Corregida.} El SD2 deja la tensión y su condición; el corte, los plazos y el tratamiento térmico pasan al SD3 (3.2.2, S-21, D-03 y RNG-15). & SD2, 2.2.3 y 2.5.1; SD3, 3.2.2 \\
\LFXfiletefila
\mbox{SD2-02} & La Figura 2.3 omite la preparación nocturna con hoja de picking y el faltante sin causa. & \textbf{Corregida.} La figura pasa a dos paneles: circuito de la hoja de picking y pérdida de la causa del faltante, distinguiendo los despachos directos. & SD2, 2.3 (proceso 3), Figura 2.3 \\
\LFXfiletefila
\mbox{SD2-03} & La Tabla 2.3 sintetiza 5 de las 12 restricciones y omite las operacionales. & \textbf{Corregida.} Se reemplaza la tabla por cuatro párrafos que cubren R-01 a R-12. & SD2, 2.5.3 \\
\LFXfiletefila
\mbox{SD2-04} & Afirmaciones que el caso no hace (WMS con soporte discontinuado; módulos del ERP abandonados). & \textbf{Corregida.} Se distingue la aplicación de preventa de 2016 sin proveedor del ERP de 2017 y se describe la cobertura real del WMS. & SD2, 2.1 \\
\LFXfiletefila
\mbox{SD2-05} & La rotación del 38 \% aparece dentro del costo de servir. & \textbf{Corregida.} Se traslada a la preparación nocturna, sin atribuirle una tasa de error no medida. & SD2, 2.2.1 y 2.3 \\
\LFXfiletefila
\mbox{SD2-06} & Cita de Boulaksil y Belkora: Casablanca en el texto y Valencia en la referencia. & \textbf{Corregida.} Se aclara que Valencia es el fabricante estudiado en Casablanca y se acota la atribución al artículo. & SD2, 2.2.4; Referencias \\
\LFXfiletefila
\mbox{SD2-07} & La Declaración de IA atribuye a la fila 2.1 la validación de la propuesta de valor. & \textbf{Corregida.} Se reemplaza por el registro de revisión del resumen ejecutivo. & SD2, Tabla 2.3 (Declaración de IA) \\
\end{longtable}
}

\Needspace{8\baselineskip}
\subsection{Subdocumento 3: Esquema de solución y alcance (T-12)}
{\linespread{1}\LFXcuerpotabla
\setlength{\parskip}{0pt}\renewcommand{\arraystretch}{1.10}\setlength{\tabcolsep}{4pt}
\begin{longtable}{F{0.09} F{0.29} F{0.40} F{0.22}}
\caption[Subdocumento 3: Esquema de solución y alcance (T-12)]{Subdocumento 3: Esquema de solución y alcance (T-12). Fuente: revisión del Informe 2 y versión corregida de LafroX.}\\[2pt]
\LFXtopruleoriginal \cab{N.º} & \cab{Observación} & \cab{Respuesta} & \cab{Sección modificada} \\ \midrule
\endfirsthead
\multicolumn{4}{@{}l@{}}{\fontsize{9pt}{11pt}\selectfont\color{lafroxGris}Tabla \thetable\ (continuación)}\\[3pt]
\LFXtopruleoriginal \cab{N.º} & \cab{Observación} & \cab{Respuesta} & \cab{Sección modificada} \\ \midrule
\endhead
\midrule
\multicolumn{4}{@{}r@{}}{\fontsize{9pt}{11pt}\selectfont\color{lafroxGris}Continúa en el folio siguiente.}\\
\endfoot
\bottomrule
\endlastfoot
\mbox{SD3-01} & Figura 3.4 ilegible; figuras delgadas (la 3.2 son cuatro cajas) y ninguna muestra qué corre en el CD, plataformas, camión y nube. & \textbf{Corregida.} Se reconstruyeron las Figuras 3.2 y 3.4, se incorporó el esquema de emplazamiento y se confirmó la revisión humana. & SD3, 3.3 y 3.4.2; Figuras 3.2, 3.4 y 3.9 \\
\LFXfiletefila
\mbox{SD3-02} & Erratas en figuras: «Gerente de operacions» y rótulos invertidos en la Figura 3.7. & \textbf{Corregida.} Se editaron las Figuras 3.4, 3.5 y 3.7 para corregir la errata y ordenar sus rótulos. & SD3, 3.4.2; Figuras 3.4, 3.5 y 3.7 \\
\LFXfiletefila
\mbox{SD3-03} & Columna «Origen» de la Tabla 3.A.1 vacía o cortada. & \textbf{Corregida.} Las 134 filas tienen origen consignado y se ajustó la tabla. & SD3-Anexos, Tabla 3.A.1; columna Origen \\
\LFXfiletefila
\mbox{SD3-04} & La prueba y el criterio de aceptación quedan fuera del T-12 (cadena del 17.1). & \textbf{Corregida.} Se incorporó al T-12 la columna de criterios de aceptación, considerando la prueba como criterio cuando corresponde. & T-12, columna Criterio de aceptación \\
\LFXfiletefila
\mbox{SD3-05} & Prioridades poco discriminantes (84 \% en Crítica o Alta). & \textbf{Corregida.} Se aplicó el techo por familia: 58 requisitos críticos, 157 altos y 55 medios; se distingue prioridad de etapa. & SD3, 3.2.3; Anexo 3.A y catálogos de requisitos \\
\LFXfiletefila
\mbox{SD3-06} & La columna «Componente» del T-12 repite la pila completa en cada fila. & \textbf{Corregida.} Se resumieron las pilas repetidas y se conservaron el módulo, la función de cada requisito y las condiciones técnicas pertinentes. & T-12, columna Componente que lo satisface \\
\LFXfiletefila
\mbox{SD3-07} & Interoperabilidad, escalabilidad, desempeño y seguridad en cuatro viñetas. & \textbf{Corregida.} Se desarrollaron los cuatro atributos, sus límites y formas previstas de verificación, con contraste de la arquitectura. & SD3, 3.4.2; referencias a los Anexos 4-D, 4-E y 4-M \\
\LFXfiletefila
\mbox{SD3-08} & Anexo 3.F reducido; fundamentos que obligan a leer SD2 y SD4. & \textbf{Corregida.} Se ampliaron los fundamentos de D-02, D-04, D-07, D-08, D-11 y D-12; las demás decisiones mantienen sus fundamentos y referencias. & SD3-Anexos, 3.F; decisiones de alcance D-01 a D-12 \\
\LFXfiletefila
\mbox{SD3-09} & Nombres de módulos distintos en SD3 y SD4 (M10 «BI» y M11 «EDI»). & \textbf{Corregida.} Los nombres coinciden con el Anexo 4-D vigente: M10 es Analítica y M11 Canal moderno. & SD3, 3.4.2, Tabla 3.4 y Figuras 3.4 y 3.7 \\
\LFXfiletefila
\mbox{SD3-10} & Meta de OTIF del mes 32 contradictoria entre el cuerpo y el Anexo 3.J. & \textbf{Corregida.} Se verificaron las metas de 90 \% en el mes 15, 93 \% en el mes 19 y más de 95 \% en el mes 32, sin redondeo. & SD3, 3.2.5; Anexo 3.J; T-17, R18-04 \\
\LFXfiletefila
\mbox{SD3-11} & Incorporación de conductores externos despachada en una frase. & \textbf{Corregida.} Se precisaron el alta conectada, la aprobación de Despacho, los permisos y la sustitución mediante suplentes previamente habilitados. & SD3, 3.4.4; Anexo 3.A, RF-06.08 \\
\LFXfiletefila
\mbox{SD3-12} & La operación no cuantifica la dotación de la mesa ni del NOC. & \textbf{Corregida.} La sección 3.4.5 indica la cobertura del NOC y las posiciones de mesa por franja, incluidos los refuerzos de septiembre y diciembre. & SD3, 3.4.5; Anexo 4-W.7 y T-15 \\
\LFXfiletefila
\mbox{SD3-13} & Revisión humana declarada sólo de forma; 3.4 declara «Ninguno» en diagramas. & \textbf{Corregida.} Todos los diagramas del Subdocumento 3 fueron elaborados y revisados por el equipo. Se actualizó la declaración de uso de IA. & SD3, declaración de uso de IA, Tabla 3.6 y diagramas de 3.4 \\
\LFXfiletefila
\mbox{SD3-14} & El SD3 no menciona ninguna de las cinco innovaciones. & \textbf{Corregida.} Se incorporaron y contrastaron las cinco innovaciones, distinguiendo la exigencia de vida útil del aporte desconectado por lote. & SD3, 3.1 y 3.4; remisiones al SD13 y T-19 \\
\end{longtable}
}

\Needspace{8\baselineskip}
\subsection{Subdocumento 4.1: Arquitectura lógica}
{\linespread{1}\LFXcuerpotabla
\setlength{\parskip}{0pt}\renewcommand{\arraystretch}{1.10}\setlength{\tabcolsep}{4pt}
\begin{longtable}{F{0.09} F{0.29} F{0.40} F{0.22}}
\caption[Subdocumento 4.1: Arquitectura lógica]{Subdocumento 4.1: Arquitectura lógica. Fuente: revisión del Informe 2 y versión corregida de LafroX.}\\[2pt]
\LFXtopruleoriginal \cab{N.º} & \cab{Observación} & \cab{Respuesta} & \cab{Sección modificada} \\ \midrule
\endfirsthead
\multicolumn{4}{@{}l@{}}{\fontsize{9pt}{11pt}\selectfont\color{lafroxGris}Tabla \thetable\ (continuación)}\\[3pt]
\LFXtopruleoriginal \cab{N.º} & \cab{Observación} & \cab{Respuesta} & \cab{Sección modificada} \\ \midrule
\endhead
\midrule
\multicolumn{4}{@{}r@{}}{\fontsize{9pt}{11pt}\selectfont\color{lafroxGris}Continúa en el folio siguiente.}\\
\endfoot
\bottomrule
\endlastfoot
\mbox{SD4.1-01} & Figura 1 (vista general) ilegible a tamaño impreso. & \textbf{Corregida.} Nueva vista vectorial de navegación por ocho capas y doce módulos, con letra mínima de 10,4 pt. & SD4, 4.1, Figura 1 \\
\LFXfiletefila
\mbox{SD4.1-02} & Matriz requerimiento → módulo → capa por grupo y no por requisito. & \textbf{Corregida.} Dieciséis obligaciones individuales con origen, componente, capa, contrato, prueba y aceptación, y cuatro recorridos críticos en el cuerpo. & SD4, Tabla 4.1-A; Anexo 4-E, Tabla A.6a \\
\LFXfiletefila
\mbox{SD4.1-03} & Leyendas de fuente de plantilla, iguales en todas las figuras. & \textbf{Corregida.} Cada figura lógica y física cita el dato que la sostiene. & SD4, leyendas de 4.1, 4.2 y 4.3; Anexos \\
\LFXfiletefila
\mbox{SD4.1-04} & Cambio de Django a Laravel sin reconocerlo frente al Informe 1. & \textbf{Corregida.} Se explica el antecedente Django, el motivo de la elección y sus efectos en contratos, pruebas y herramientas. & SD4, 4.1.7.1; ADR-23 \\
\LFXfiletefila
\mbox{SD4.1-05} & Versiones recientes (Laravel 13, PHP 8.5, Angular 22, TypeScript 6.0) por demostrar. & \textbf{Corregida.} Tabla de lanzamiento, restricciones y soporte de ocho productos con fuentes oficiales y política de actualización. & SD4, Anexo 4-P, Tabla A.20 \\
\LFXfiletefila
\mbox{SD4.1-06} & Referencia a la «sección 2.5» inexistente del caso. & \textbf{Corregida.} Se cita la clasificación correcta (RT-06.01). & SD4, 4.3.1 \\
\LFXfiletefila
\mbox{SD4.1-07} & Párrafo con redacción defectuosa sobre disponibilidad y redundancia. & \textbf{Corregida.} Se reescribe el párrafo. & SD4, 4.3.1 \\
\LFXfiletefila
\mbox{SD4.1-08} & Credencial de turno de 14 h frente a rutas de hasta 14,5 h; manifiesto de 26 h. & \textbf{Corregida.} Asignación emitida dentro de la hora previa a salir, válida hasta el fin programado más 30 min y con máximo de 16 h; manifiesto con vigencia remanente mínima de 25 h y dos relevos. & SD4, 4.1 (identidad y conductores externos); ADR-06; Anexo 4-M \\
\LFXfiletefila
\mbox{SD4.1-09} & Declaración de IA «Alto» en texto y diagramas con un solo revisor; 4.3 declara «Ninguno». & \textbf{Corregida.} Cinco revisores se reparten las secciones, anexos y el T-11, cada uno con el contenido que verificó. & SD4, Tabla 29 (Declaración de IA) \\
\end{longtable}
}

\Needspace{8\baselineskip}
\subsection{Subdocumento 4.2 y 4.3: Arquitectura física y centros de datos (T-11)}
{\linespread{1}\LFXcuerpotabla
\setlength{\parskip}{0pt}\renewcommand{\arraystretch}{1.10}\setlength{\tabcolsep}{4pt}
\begin{longtable}{F{0.09} F{0.29} F{0.40} F{0.22}}
\caption[Subdocumento 4.2 y 4.3: Arquitectura física y centros de datos (T-11)]{Subdocumento 4.2 y 4.3: Arquitectura física y centros de datos (T-11). Fuente: revisión del Informe 2 y versión corregida de LafroX.}\\[2pt]
\LFXtopruleoriginal \cab{N.º} & \cab{Observación} & \cab{Respuesta} & \cab{Sección modificada} \\ \midrule
\endfirsthead
\multicolumn{4}{@{}l@{}}{\fontsize{9pt}{11pt}\selectfont\color{lafroxGris}Tabla \thetable\ (continuación)}\\[3pt]
\LFXtopruleoriginal \cab{N.º} & \cab{Observación} & \cab{Respuesta} & \cab{Sección modificada} \\ \midrule
\endhead
\midrule
\multicolumn{4}{@{}r@{}}{\fontsize{9pt}{11pt}\selectfont\color{lafroxGris}Continúa en el folio siguiente.}\\
\endfoot
\bottomrule
\endlastfoot
\mbox{SD4.2-01} & La sala no declara superficie ni tiene plano acotado (RT-06.03); recintos de más para dos racks. & \textbf{Corregida.} Se declara la superficie útil del recinto y la Figura 29 pasa a plano acotado por zonas. & SD4, 4.3.1, tabla de superficie y Figura 29 \\
\LFXfiletefila
\mbox{SD4.2-02} & Sin alternativa si el CLIENTE no aprueba la residencia en sa-east-1 y us-east-1. & \textbf{Pendiente.} El ADR-19 mantiene la aprobación del CLIENTE como hito de la Etapa 1, sin arquitectura alternativa. Se corrige para la Propuesta Final. & SD4, Anexo 4-O, ADR-19 \\
\LFXfiletefila
\mbox{SD4.2-03} & Estudio de cobertura dentro de las cámaras diferido. & \textbf{Corregida.} Estudio con diseño predictivo por zona y verificación en sitio en el mes 2 (paquete 2.2.2). & SD4, Anexo 4-W.13; T-14, 2.2.2 \\
\LFXfiletefila
\mbox{SD4.2-04} & Figura 18 (topología AWS) con rótulos de 6 a 7 pt. & \textbf{Corregida.} Figura redibujada a 9 pt. & SD4, 4.2, Figura 18 \\
\LFXfiletefila
\mbox{SD4.2-05} & La Tabla 14 no explica el total por máximo coincidente. & \textbf{Corregida.} Nota que explica que el total es el máximo agregado coincidente y no la suma de máximos. & SD4, 4.2.6, Tabla 14 \\
\LFXfiletefila
\mbox{SD4.2-06} & LafroX provee la cadena de frío, contra la exclusión del Cap. 11 del caso. & \textbf{Corregida.} El CLIENTE adquiere los sensores y kits según la especificación del T-11; LafroX los instala, integra y mantiene. Alineado en SD3, SD6, T-9 y SD7. & SD4, 4.2.3; T-11; SD6, 6.1; SD7, 5.1.4 \\
\LFXfiletefila
\mbox{SD4.2-07} & El T-11 no declara ciclo de vida ni plan de reposición por equipo. & \textbf{Corregida.} Tabla de ciclo de vida, repuestos y reposición por dispositivo en los 56 meses. & T-11, ciclo de vida por dispositivo \\
\LFXfiletefila
\mbox{SD4.2-08} & Sobra la frase «los niveles del 7.2 son un medio, no un fin». & \textbf{Corregida.} Frase eliminada. & SD4, 4.3 \\
\LFXfiletefila
\mbox{SD4.2-09} & Riesgo regulatorio y de latencia de la residencia fuera de Chile no llevado al plan de riesgos. & \textbf{Parcial.} El T-16 registra la falta de aprobación de regiones de nube como riesgo vigente; la latencia no se cuantifica. Lo que falta se completa para la Propuesta Final. & T-16; SD8, Anexo 8.F \\
\end{longtable}
}

\Needspace{8\baselineskip}
\subsection{Subdocumento 5: Modelo y gestión de datos}
{\linespread{1}\LFXcuerpotabla
\setlength{\parskip}{0pt}\renewcommand{\arraystretch}{1.10}\setlength{\tabcolsep}{4pt}
\begin{longtable}{F{0.09} F{0.29} F{0.40} F{0.22}}
\caption[Subdocumento 5: Modelo y gestión de datos]{Subdocumento 5: Modelo y gestión de datos. Fuente: revisión del Informe 2 y versión corregida de LafroX.}\\[2pt]
\LFXtopruleoriginal \cab{N.º} & \cab{Observación} & \cab{Respuesta} & \cab{Sección modificada} \\ \midrule
\endfirsthead
\multicolumn{4}{@{}l@{}}{\fontsize{9pt}{11pt}\selectfont\color{lafroxGris}Tabla \thetable\ (continuación)}\\[3pt]
\LFXtopruleoriginal \cab{N.º} & \cab{Observación} & \cab{Respuesta} & \cab{Sección modificada} \\ \midrule
\endhead
\midrule
\multicolumn{4}{@{}r@{}}{\fontsize{9pt}{11pt}\selectfont\color{lafroxGris}Continúa en el folio siguiente.}\\
\endfoot
\bottomrule
\endlastfoot
\mbox{SD5-01} & Figura 5.1 ilegible (texto de 1 a 2 pt). & \textbf{Corregida.} Se reemplaza por una vista de dominios y relaciones principales; el detalle queda en las Figuras 5.2 a 5.10. & SD5, 5.1, Figura 5.1 \\
\LFXfiletefila
\mbox{SD5-02} & Vistas por dominio cerca de 8 pt y título repetido en la imagen. & \textbf{Corregida.} Se retiraron los títulos duplicados de las imágenes. El texto interno de las nueve figuras mide entre 10,68 y 11,66 pt en el PDF. & SD5, Figuras 5.2 a 5.10 \\
\LFXfiletefila
\mbox{SD5-03} & Contradicción sobre «M» entre el cuerpo y el glosario. & \textbf{Corregida.} «M» queda reservada a módulos; los plazos se escriben «mes N». & SD5, introducción; Anexo 5-A, glosario \\
\LFXfiletefila
\mbox{SD5-04} & No se declara de dónde sale el umbral de 0,3 \% de desviación de inventario. & \textbf{Corregida.} Se fundamenta como control de riesgo de LafroX para la carga inicial. & SD5, 5.3 \\
\LFXfiletefila
\mbox{SD5-05} & El SD5 no muestra el cálculo de los volúmenes a migrar. & \textbf{Corregida.} Cálculo por conjunto (registros × meses × KB): 16.054.780 KB de origen y 32,1 GB de destino, con sensibilidad. & SD5, 5.3 \\
\LFXfiletefila
\mbox{SD5-06} & Prosa telegráfica, cargada de barras, siglas y remisiones. & \textbf{Corregida.} Se amplió la redacción explicativa a 5.2.5 y 5.2.6, con las fórmulas, los umbrales y los plazos de los criterios técnicos. & SD5, introducción y 5.1 a 5.4 \\
\LFXfiletefila
\mbox{SD5-07} & Dos nomenclaturas para los subdocumentos (SD1–SD5 y SD1–SD5) y nombre del archivo en la portada. & \textbf{Corregida.} Se usa sólo «SD» y se retira el nombre del archivo de la portada. & SD5, cuerpo y portada \\
\LFXfiletefila
\mbox{SD5-08} & Declaración de IA «Alto» con un solo revisor. & \textbf{Corregida.} La declaración identifica a Bastian Trejo y Alex Aravena como revisores y detalla qué verificaron en cada una de las 16 filas. & SD5, Declaración de IA, Tabla 5.4 \\
\LFXfiletefila
\mbox{SD5-09} & Los relevos de 8/16 h y la credencial de turno no se explican en el SD5. & \textbf{Corregida.} Se explica la función de cada plazo: identidad de consulta de 24 h, credencial del turno y manifiesto de 26 h para relevos inscritos, con su prueba de corte y cambio de turno. & SD5, 5.2, Tabla 5.1 \\
\end{longtable}
}

\Needspace{8\baselineskip}
\subsection{Subdocumento 6: Metodologías (T-9 y T-10)}
{\linespread{1}\LFXcuerpotabla
\setlength{\parskip}{0pt}\renewcommand{\arraystretch}{1.10}\setlength{\tabcolsep}{4pt}
\begin{longtable}{F{0.09} F{0.29} F{0.40} F{0.22}}
\caption[Subdocumento 6: Metodologías (T-9 y T-10)]{Subdocumento 6: Metodologías (T-9 y T-10). Fuente: revisión del Informe 2 y versión corregida de LafroX.}\\[2pt]
\LFXtopruleoriginal \cab{N.º} & \cab{Observación} & \cab{Respuesta} & \cab{Sección modificada} \\ \midrule
\endfirsthead
\multicolumn{4}{@{}l@{}}{\fontsize{9pt}{11pt}\selectfont\color{lafroxGris}Tabla \thetable\ (continuación)}\\[3pt]
\LFXtopruleoriginal \cab{N.º} & \cab{Observación} & \cab{Respuesta} & \cab{Sección modificada} \\ \midrule
\endhead
\midrule
\multicolumn{4}{@{}r@{}}{\fontsize{9pt}{11pt}\selectfont\color{lafroxGris}Continúa en el folio siguiente.}\\
\endfoot
\bottomrule
\endlastfoot
\mbox{SD6-01} & Regla de valor ganado truncada («una desviación mayor al 10»). & \textbf{Corregida.} Regla completa en el cuerpo y en el T-9, con tabla de umbrales: SPI o CPI bajo 0,90 o desviación de más de 10 \% en la ruta crítica obliga a un plan de recuperación en cinco días hábiles, aprobado por el Comité de Proyecto. & SD6, 6.1.5; T-9, 1.5, Tabla T-9.4 \\
\LFXfiletefila
\mbox{SD6-02} & PMBOK «adaptado» sin explicar la adaptación ni la convivencia de RUP, Kanban e iteraciones con los hitos. & \textbf{Corregida.} La gestión distingue los hitos obligatorios de la programación de actividades y recursos. Cada etapa sigue una cascada, con frentes paralelos en la construcción. & SD6, 6.1 y 6.2; T-9, 1.1. \\
\LFXfiletefila
\mbox{SD6-03} & Control de cambios sin las reglas del Art. 72°. & \textbf{Corregida.} Plazo de análisis de impacto (5 o 10 días hábiles), tope acumulado con alerta al 15 \% y límite del 20 \%, y regla para cambios en una etapa con fecha fija. & SD6, 6.1.4; T-9, 1.4, Tabla T-9.3 \\
\LFXfiletefila
\mbox{SD6-04} & Escalamiento sin destinatario, plazo ni disparador. & \textbf{Corregida.} Tres niveles con disparadores observables y plazos (mismo día, 2 y 5 días hábiles), y escalamiento inmediato ante amenazas a seguridad, continuidad o hitos; coincide con el SD8. & SD6, 6.1.6, Figura 6.2; T-9, 1.6, Tabla T-9.5 \\
\LFXfiletefila
\mbox{SD6-05} & Matriz de comunicaciones genérica (sin terreno, turno de noche, terceros ni planificador). & \textbf{Corregida.} Matriz con canales de terreno: inicio de jornada en andén, cambio de turno a las 22:00, conductor externo al asignarse, aviso en el terminal, planificador y sindicato. & SD6, 6.1.2; T-9, 1.2, Tabla T-9.2 \\
\LFXfiletefila
\mbox{SD6-06} & Registro de 17 interesados sin influencia ni interés, frente a 19 en SD2 y SD3. & \textbf{Corregida.} Registro de los 19 actores con la influencia e interés del SD2, estrategia por celda, momento, responsable y canal, alineado con el SD3. & SD6-Anexos, Tabla 6.A.1 \\
\LFXfiletefila
\mbox{SD6-07} & Párrafos de evaluadores externos repetidos en SD6 y T-9. & \textbf{Corregida.} Se retiran del SD6 y del T-9; el dimensionamiento queda en el plan de calidad. & SD6, 6.1.3; T-9, 1.3 \\
\LFXfiletefila
\mbox{SD6-08} & Sin método para nivelar la capacidad en los solapamientos 13–15 y 19–20. & \textbf{Corregida.} Control semanal de horas asignadas frente a la curva de plazas del T-15 y el paquete 1.2.4, con reasignación o reserva protegida y escalamiento. & SD6, 6.1.5; T-9, 1.5 \\
\LFXfiletefila
\mbox{SD6-09} & RUP justificado en una frase, sin comparar alternativas. & \textbf{Corregida.} Se elige la cascada por etapa debido a los hitos de aceptación, las interfaces de terceros y la capacidad del área de TI. Los referentes validan las reglas de terreno antes de construir. & SD6, 6.2; T-10, introducción. \\
\LFXfiletefila
\mbox{SD6-10} & Requerimientos, arquitectura evolutiva, refactorización y deuda técnica en tres líneas. & \textbf{Corregida.} Gestión de requerimientos (casos de uso, refinamiento con el referente del CLIENTE, trazabilidad y cambios) y tratamiento de la deuda por clase con plazos. & SD6, 6.2.1; T-10, secciones 1 y 2, Tabla T-10.1 \\
\LFXfiletefila
\mbox{SD6-11} & No dice cómo se desarrolla y prueba lo propio de Puelche (desconexión, −22 °C, ERP sin documentar, cadenas, planificador). & \textbf{Corregida.} Seis condiciones del caso con su forma de desarrollo y su prueba antes de la marcha blanca, más la transferencia al área de TI de cuatro personas y el tiempo de salida al mercado. & SD6, 6.2 y 6.2.4, Tabla 6.3; T-10, sección 5 \\
\LFXfiletefila
\mbox{SD6-12} & El T-10 habla de «los cinco ambientes» sin nombrarlos ni definir promoción y aprobación. & \textbf{Corregida.} Tabla de los cinco ambientes con lo que se despliega, la puerta de entrada y quién aprueba. & T-10, Tabla T-10.2; SD6, Figura 6.4 \\
\LFXfiletefila
\mbox{SD6-13} & Iteraciones, reuniones y liberaciones sin fecha en la planificación. & \textbf{Corregida.} En la cascada por etapa, las revisiones se vinculan a las condiciones de salida de cada fase, al seguimiento semanal, a cada despliegue en QA y a la certificación. Las fechas de salida de fase y certificación corresponden a los hitos del T-15. & SD6, 6.2.3; T-10, Revisiones; T-15, 6 \\
\LFXfiletefila
\mbox{SD6-14} & Cuerpo de ocho páginas con una sola figura, cuya fuente es equivocada; sin diagramas de gobierno, iteración ni pipeline. & \textbf{Corregida.} Cuerpo de 18 páginas con cuatro figuras (modelo en tres niveles, escalamiento, control de cambios y cadena de entrega) y fuentes corregidas. & SD6, Figuras 6.1 a 6.4 \\
\LFXfiletefila
\mbox{SD6-15} & T-9 y T-10 repiten el SD6 en lugar de detallar. & \textbf{Corregida.} Los formularios detallan la adaptación metodológica, el control de cambios, los umbrales, el escalamiento, el contenido de las actas, la deuda técnica, los ambientes y las revisiones. & T-9; T-10 \\
\LFXfiletefila
\mbox{SD6-16} & El SD7 remite a las iteraciones de la 6.2, pero el cuerpo del SD6 no fija su duración (sólo el T-10). & \textbf{Corregida.} El SD6 y el T-10 describen la cascada por etapa. La programación del T-15 se organiza por fases y paquetes de trabajo, con frentes paralelos en construcción. & SD6, 6.2 y 6.2.3; T-10; T-15, 6 \\
\LFXfiletefila
\mbox{SD6-17} & Declaración de IA «Alto» en todas las secciones; la revisión declarada es de coherencia, no de autoría. & \textbf{Corregida.} Revisores distribuidos por sección, cada uno con el contenido que verificó. & SD6, Declaración de IA \\
\end{longtable}
}

\Needspace{8\baselineskip}
\subsection{Subdocumento 7: Plan de trabajo, EDT y cronograma (T-14, T-15 y T-18)}
{\linespread{1}\LFXcuerpotabla
\setlength{\parskip}{0pt}\renewcommand{\arraystretch}{1.10}\setlength{\tabcolsep}{4pt}
\begin{longtable}{F{0.09} F{0.29} F{0.40} F{0.22}}
\caption[Subdocumento 7: Plan de trabajo, EDT y cronograma (T-14, T-15 y T-18)]{Subdocumento 7: Plan de trabajo, EDT y cronograma (T-14, T-15 y T-18). Fuente: revisión del Informe 2 y versión corregida de LafroX.}\\[2pt]
\LFXtopruleoriginal \cab{N.º} & \cab{Observación} & \cab{Respuesta} & \cab{Sección modificada} \\ \midrule
\endfirsthead
\multicolumn{4}{@{}l@{}}{\fontsize{9pt}{11pt}\selectfont\color{lafroxGris}Tabla \thetable\ (continuación)}\\[3pt]
\LFXtopruleoriginal \cab{N.º} & \cab{Observación} & \cab{Respuesta} & \cab{Sección modificada} \\ \midrule
\endhead
\midrule
\multicolumn{4}{@{}r@{}}{\fontsize{9pt}{11pt}\selectfont\color{lafroxGris}Continúa en el folio siguiente.}\\
\endfoot
\bottomrule
\endlastfoot
\mbox{SD7-01} & Primer nivel de la EDT organizado por fases RUP y no por entregables. & \textbf{Corregida.} Los diez productos forman el primer nivel de la EDT. El T-14 presenta la vista general, una figura por producto y el diccionario de las 60 cuentas y los 230 paquetes. & SD7, 7.1, Figura 7.1; T-14, secciones 1 y 2 \\
\LFXfiletefila
\mbox{SD7-02} & Las etapas aparecen como ramas separadas de la EDT. & \textbf{Corregida.} Las figuras de la EDT agrupan los paquetes por producto, sin ramas por etapa. El diccionario registra la etapa en el campo «Período» de cada paquete. & SD7, 7.1; T-14, secciones 1 y 2 \\
\LFXfiletefila
\mbox{SD7-03} & Estimación por clases fijas (todo módulo vale 960 HH). & \textbf{Corregida.} Estimación por tamaño de cada módulo: flujos S/M/C, interfaces y dispositivos trazables al T-12, con actividades de diseño, construcción, pruebas e integración. & T-15, 4.3 (estimación por tamaño) \\
\LFXfiletefila
\mbox{SD7-04} & PERT degenerado (O = 0,75M y P = 1,25M para todos). & \textbf{Corregida.} O y P con fundamento propio por módulo, tipo de paquete y riesgo; la banda del 25 \% queda sólo como sensibilidad. & T-15, 4.3 y 5 \\
\LFXfiletefila
\mbox{SD7-05} & Capacidad de calidad insuficiente (faltan 217,60 HH) e implantación sin división que la provea. & \textbf{Corregida.} Tres refuerzos de pruebas (1.081,60 HH frente a 1.049,60), área de Implantación con 9 personas y 18 asistentes por obra; 137 plazas conciliadas con el SD1, con la mesa y el SOC aparte de su planta técnica. & T-15, 5.7 y 6.3; SD1, 1.2 \\
\LFXfiletefila
\mbox{SD7-06} & Once módulos de la Etapa 1 en paralelo, sobre el techo de desarrolladores del SD1. & \textbf{Corregida.} Los módulos de la Etapa 1 se construyen en tres olas, con equipos de 3 a 7 personas y un máximo de 31 personas simultáneas en módulos. Desarrollo llega a 34 plazas en el mes 7 (una de dirección y 33 de construcción), dentro de las 36 plazas previstas para el grupo y de los 48 desarrolladores de planta. & T-15, 4.4, 5.7 y 6.1 \\
\LFXfiletefila
\mbox{SD7-07} & Desajuste entre el plazo de revisión del acta (Art. 18.3) y el paso a producción. & \textbf{Parcial.} Tras diez días hábiles de revisión, el acta de H7 se firma el 16-05-2028 y la producción comienza el 17-05-2028. Para H12, las fechas son el 17-10-2028 y el 18-10-2028, respectivamente. La carta Gantt se rehace con esas fechas. Lo que falta se completa para la Propuesta Final. & T-15, Tabla 14 y sección 5.5; SD7, 7.3.1, Figura 7.5 \\
\LFXfiletefila
\mbox{SD7-08} & Los once módulos de la Etapa 1 arrancan la misma semana; la secuencia por dependencias que describe el texto no se modela como orden de construcción. & \textbf{Parcial.} La ola A comienza el 10-06-2027 y la B el 18-06-2027, tras aprobar los contratos de inventario. En la ola C, cobranza (4.3.5) parte el 11-08-2027 en comienzo a comienzo con reparto (4.3.3); los demás módulos parten al terminar su precedente. La carta Gantt se rehace con esa programación. Lo que falta se completa para la Propuesta Final. & T-15, 5.1; SD7, 7.3.2 \\
\LFXfiletefila
\mbox{SD7-09} & Red incompleta: 102 de 168 paquetes sin predecesora (migración, cadenas, innovaciones). & \textbf{Corregida.} La red reúne 408 eventos y 747 relaciones. Sólo 1.1.1 y 2.4.1 parten sin predecesora al inicio del contrato. El paquete 8.1.4 depende de la especificación del SD4 y de la ventana de homologación móvil. & T-15, 5.1 y 6.1; SD7, Anexo 7.B \\
\LFXfiletefila
\mbox{SD7-10} & CPM calculado sobre fechas publicadas y no sobre duraciones. & \textbf{Corregida.} El CPM se calcula a partir de las duraciones y las predecesoras de la red por paquetes. & T-15, 5 \\
\LFXfiletefila
\mbox{SD7-11} & Iteraciones de dos semanas sólo como plantilla; no se ven en la Gantt de construcción. & \textbf{Corregida.} Cada etapa sigue una cascada: la construcción avanza en frentes paralelos, sin iteraciones. La revisión quincenal controla el avance de los paquetes. & SD7, 7.2 y 7.3.8; T-15, 6 \\
\LFXfiletefila
\mbox{SD7-12} & Pruebas concentradas después de construir (cierre en V). & \textbf{Corregida.} Cada cambio pasa por pruebas unitarias y de contrato. En cada despliegue a QA se ejecutan pruebas de integración; las pruebas ampliadas se realizan durante la certificación. & SD7, 7.2; T-15, 6; SD9, 9.3 y Anexo 9.E.1 \\
\LFXfiletefila
\mbox{SD7-13} & Gantt con rótulos de 5 a 6 pt, encabezados superpuestos y códigos sin definir. & \textbf{Parcial.} Las figuras de la EDT, su cobertura, la red de precedencias y la ruta de H12 son vectoriales y sus rótulos usan letra de al menos 9 pt. La carta Gantt de 56 meses y sus detalles se rehacen con los códigos vigentes. Lo que falta se completa para la Propuesta Final. & SD7, Figuras 7.1 a 7.3 y 7.5; Gantt, Figuras 7.7 a 7.14 \\
\LFXfiletefila
\mbox{SD7-14} & Filas agrupadas que mezclan paquetes de cuentas distintas. & \textbf{Parcial.} La EDT y el diccionario agrupan los 230 paquetes por producto y cuenta de control. La carta Gantt y sus láminas de detalle se rehacen con esa misma agrupación. Lo que falta se completa para la Propuesta Final. & T-14, secciones 1 y 2; SD7, 7.3.2 \\
\LFXfiletefila
\mbox{SD7-15} & Defecto crítico durante el congelamiento de septiembre en la marcha blanca de la Etapa 2. & \textbf{Corregida.} Del 1 al 25 de septiembre se registra la severidad, se contiene el defecto y se escala a la Contraparte Técnica sin alterar Producción. Del 26 al 30 sólo se despliega un cambio de emergencia con autorización expresa, después de probarlo en Preproducción y fuera de la ventana de despacho. & SD7, 7.3.5; T-18, 6.4 \\
\LFXfiletefila
\mbox{SD7-16} & Reversión: 10 min en la Tabla 7.8 y 40 min en el 7.3.8 sin explicación. & \textbf{Corregida.} Se explican los niveles: max(10, 30) + 10 = 40 min, con decisión a las 04:45 y restauración a las 05:25. & SD7, 7.3, Tabla 7.8 \\
\LFXfiletefila
\mbox{SD7-17} & El intercambio con la principal cadena parte en el mes 13, sin margen ante agosto de 2028. & \textbf{Parcial.} El T-14 y el T-15 sitúan el inicio del intercambio con la principal cadena (5.3.1) el 01-10-2027, en el mes 9. La carta Gantt se rehace con esa fecha. Lo que falta se completa para la Propuesta Final. & T-14, sección 2, paquete 5.3.1; T-15, 5.1 y 6.1 \\
\LFXfiletefila
\mbox{SD7-18} & Sin evidencia al proveedor de lácteos entre el inicio y el mes 12. & \textbf{Corregida.} El acta de inicio 1.1.3 incorpora la evidencia para el proveedor de lácteos: el CLIENTE valida el formato y el canal de un expediente inicial con los registros de recepción y lote, el recorrido hasta el despacho y los destinatarios recuperables. Hasta habilitar la interfaz 5.2.1, Calidad mantiene la evidencia con los registros existentes. & T-14, 1.1.3 y 5.2.1 \\
\LFXfiletefila
\mbox{SD7-19} & Errata en la línea base de alcance de la Etapa 1: figura «Mes 1 (H1)», pero la línea base se firma en el mes 2. & \textbf{Corregida.} La línea base de la Etapa 1 (2.1.1) se elabora en el mes 1 y se aprueba en H1 (mes 2). La línea base de la Etapa 2 (2.1.2) distingue también la elaboración de la aprobación en H8. & T-14, diccionario \\
\LFXfiletefila
\mbox{SD7-20} & Tablas 7.2 y 7.8 con párrafos en las celdas. & \textbf{Corregida.} Celdas breves y explicación en prosa contigua. & SD7, Tablas 7.2 y 7.8 \\
\LFXfiletefila
\mbox{SD7-21} & El 7.3.8 arrastra contenido ajeno al plan, duplicado en el T-18. & \textbf{Corregida.} El apartado 7.3.8 conserva la estimación, la programación por fases, el calendario y la implantación. La regla de precio, la coordinación de la reserva de stock y los objetivos de recuperación se retiran de ese apartado. Las condiciones técnicas de las pruebas se detallan en el T-18, sección 6.5. & SD7, 7.3.8; T-18, 6.5 \\
\LFXfiletefila
\mbox{SD7-22} & Siglas SD4, SD8, PT/H y RC sin definir; Figura 7.20 remite a un detalle inexistente. & \textbf{Corregida.} El SD7 define las siglas SD4 y SD8, y sus leyendas distinguen RC (ruta crítica) de PT/H (los demás paquetes e hitos). La Figura 7.3 presenta la red; el detalle de las dependencias está en el Anexo 7.B. & SD7, introducción, 7.2.1, Figura 7.3 y Anexo 7.B \\
\LFXfiletefila
\mbox{SD7-23} & Volumen: 66 páginas de cuerpo y 250 de formularios y anexos. & \textbf{Parcial.} El cuerpo del SD7 pasó de 116 a 90 páginas. Los formularios y anexos pasaron de 376 a 324 páginas: 142 del T-14, 132 del T-15, 24 del T-18 y 26 de los Anexos del SD7. Lo que falta se completa para la Propuesta Final. & SD7; T-14, T-15, T-18 y Anexos del SD7 \\
\LFXfiletefila
\mbox{SD7-24} & Redacción defensiva («no acredita», «no demuestra», «no se presume»). & \textbf{Corregida.} Se retiraron las fórmulas defensivas del SD7, el T-14, el T-15 y el T-18. Se conservan sólo las negaciones que fijan una condición técnica. & SD7; T-14; T-15; T-18 \\
\LFXfiletefila
\mbox{SD7-25} & Falta la ruta crítica de la construcción de la Etapa 1, no sólo la del tramo final. & \textbf{Parcial.} En la construcción de la Etapa 1, 4.3.3 y 4.3.5 tienen 3 días hábiles de holgura; 4.3.1 y 6.1.1 tienen 7. Los paquetes más ajustados de infraestructura y nube tienen 1 día. La Figura 7.5 muestra la ruta de H12. La vista crítica de la Gantt se rehace. Lo que falta se completa para la Propuesta Final. & SD7, 7.3.1, Figuras 7.5 y 7.6; T-15, sección 2 \\
\LFXfiletefila
\mbox{SD7-26} & Sólo las 119 actividades de los paquetes que terminan antes del H2 están detalladas; el resto queda a planificación gradual. & \textbf{Parcial.} Hasta H2 se detallan 119 actividades de 60 paquetes. Entre H2 y H5 se programan 385 actividades de 76 paquetes, por plantilla y quincena. Los 35 paquetes posteriores a H5 quedan en planificación gradual, con sus horizontes en la Figura 3. Lo que falta se completa para la Propuesta Final. & T-15, 6.1, Figura 3 \\
\end{longtable}
}

\Needspace{8\baselineskip}
\subsection{Subdocumento 8: Plan de riesgos (T-16)}
{\linespread{1}\LFXcuerpotabla
\setlength{\parskip}{0pt}\renewcommand{\arraystretch}{1.10}\setlength{\tabcolsep}{4pt}
\begin{longtable}{F{0.09} F{0.29} F{0.40} F{0.22}}
\caption[Subdocumento 8: Plan de riesgos (T-16)]{Subdocumento 8: Plan de riesgos (T-16). Fuente: revisión del Informe 2 y versión corregida de LafroX.}\\[2pt]
\LFXtopruleoriginal \cab{N.º} & \cab{Observación} & \cab{Respuesta} & \cab{Sección modificada} \\ \midrule
\endfirsthead
\multicolumn{4}{@{}l@{}}{\fontsize{9pt}{11pt}\selectfont\color{lafroxGris}Tabla \thetable\ (continuación)}\\[3pt]
\LFXtopruleoriginal \cab{N.º} & \cab{Observación} & \cab{Respuesta} & \cab{Sección modificada} \\ \midrule
\endhead
\midrule
\multicolumn{4}{@{}r@{}}{\fontsize{9pt}{11pt}\selectfont\color{lafroxGris}Continúa en el folio siguiente.}\\
\endfoot
\bottomrule
\endlastfoot
\mbox{SD8-01} & 13 riesgos internos del proveedor declarados al CLIENTE. & \textbf{Parcial.} El T-16 reúne 36 riesgos y 8.E.1 define 15 controles CT-01 a CT-15, con responsables y evidencia de cierre. CT-03 integra identidad conectada/local y CT-15 conserva stock. ERP, stock e identidad se sincronizan documentalmente con SD4/SD9; faltan implementación, eficacia y las restantes dependencias de diseño/prueba de SD4, SD9 y SD13. Lo que falta se completa para la Propuesta Final. & SD8, 8.2.1; SD8-Anexos, 8.D.3 y 8.E.1; T-16 \\
\LFXfiletefila
\mbox{SD8-02} & Riesgos del CLIENTE sin contraparte, pedido ni fecha. & \textbf{Corregida.} Cada riesgo nombra la contraparte, qué se pide, el plazo y la contingencia si no responde. & SD8-Anexos, 8.A.2 y 8.F.1; T-16 \\
\LFXfiletefila
\mbox{SD8-03} & No hay riesgo residual. & \textbf{Corregida.} Se proyecta la valoración posterior de los 36 riesgos: 1 crítico, 4 altos, 20 medios y 11 bajos. p1 depende de ejecutar el tratamiento completo y verificar evidencia; sin ella se conserva la valoración inicial. No se acredita eficacia observada. & SD8, 8.3.1; SD8-Anexos, 8.F.2; T-16, resumen residual \\
\LFXfiletefila
\mbox{SD8-04} & Riesgos en zona Alta o Crítica sin plan adicional. & \textbf{Corregida.} Ajuste documental completado. R-002 permanece Crítica; R-006, R-029, R-030 y R-031, Altos. Se completan las cinco fichas y la matriz 8.F.2a con prueba, evidencia, verificador, disparador, plazo y decisión. Sólo evidencia satisfactoria permite proponer reevaluación; resultado insuficiente o fallido mantiene el bloqueo y exige corregir/repetir. Se distingue valoración inicial, objetivo condicionado y eficacia acreditada; se revisa suficiencia de HH sin duplicaciones. Falta acreditar eficacia y nivel manejable. & SD8, 8.3.1; SD8-Anexos, 8.F.1, 8.F.2 y 8.F.2a; 8.G.1; T-16 \\
\LFXfiletefila
\mbox{SD8-05} & Faltan riesgos del Cap. 19 (sindicato, rotación, interfaces ERP/WMS, Concepción, peak) y obsolescencia, bloqueo y escalabilidad. & \textbf{Corregida.} Ajuste documental completado. La matriz 8.D.2 coteja los diez temas: R-014, R-017 y R-028 a R-035, más R-036 por divergencia del perfil real de carga. R-036 es Medio (p0=p1=30 \%, I=4, E=0,12), sin descuento por eficacia futura. CT-14 conserva el peak conocido, crecimiento comprometido y corrección del diseño a cargo de LafroX. Se actualizan RBS, análisis y reservas para 36 riesgos, sin duplicar HH. Cobertura documental completada; faltan validar el perfil, ejecutar pruebas y acreditar eficacia. & SD8-Anexos, 8.D.2, 8.B.1–8.B.2, 8.E.1 (CT-14), 8.F y 8.H; T-16, R-036; SD8, Figura 8.1 \\
\LFXfiletefila
\mbox{SD8-06} & Puntos únicos de falla del SD4 no llevados al registro. & \textbf{Corregida.} Ajuste documental completado. R-034 conserva la dependencia del ERP y CT-15 el stock con retención local durable y autoridad única. CT-03 incorpora identidad conectada/local conforme a SD4/SD9, sin renumerar controles. La matriz 8.D.3 enlaza diseño, protocolos, campañas, responsables y evidencia; SD4 y T-13/T-17 incorporan referencias recíprocas. Se conservan los 36 riesgos y sus valoraciones. Sincronización documental completada; faltan evidencia de implementación, resultados de ejecución y actas de aceptación. & SD8-Anexos, 8.D.3 y 8.E.1; T-16, R-034; SD4, Anexos 4-M/4-V (Tabla 4-V.1); SD9, T-13, 7.5 y T-17, 3.1 \\
\LFXfiletefila
\mbox{SD8-07} & No se aplica FMEA, árbol de fallas ni simulación. & \textbf{Corregida.} FMEA de 15 modos con S/O/D, NPR, controles y puntuación posterior. & SD8, 8.2.3; Anexo 8.G.1 \\
\LFXfiletefila
\mbox{SD8-08} & Sin análisis costo-beneficio. & \textbf{Corregida.} Quince cálculos en HH con sensibilidad y justificación de los retornos menores que 1. & SD8-Anexos, 8.F.1 y 8.F.3 \\
\LFXfiletefila
\mbox{SD8-09} & Acciones sin plazo, esfuerzo ni contingencia. & \textbf{Corregida.} Las 36 fichas fijan responsables, perfiles, HH, plazo preventivo, actuación contingente y evidencia. La distribución mensual por perfil se propone en 8.H.3; requiere nivelación y autorización, no acredita capacidad comprometida ni ejecución. & SD8-Anexos, 8.F.1 y 8.H.2–8.H.3 \\
\LFXfiletefila
\mbox{SD8-10} & Reservas sin derivar ni reflejo en el cronograma; 3.072 HH del SD7 no conciliadas. & \textbf{Corregida.} La Propuesta Final aprueba 160 HH de gestión, asignadas por titular y turno en cuatro ciclos de 40 HH (meses 12, 18, 24 y 40), con cobertura de PMO para JP y de DevOps para SRE. La exposición inicial es 1.199,52 HH y el total posterior 1.017,53 HH. Las 3.072 HH de protección exclusiva de E1 integran las 203.548,62 HH programadas, con crédito inicial cero. & SD8, 8.3.2; SD8-Anexos, 8.B.1 y 8.H \\
\LFXfiletefila
\mbox{SD8-11} & Escala de impacto genérica, factores sin definir, bandas solapadas en 70 \%, frase que niega los umbrales de zona y dos probabilidades por riesgo. & \textbf{Corregida.} Factores, bandas y umbrales únicos definidos en el cuerpo; se retira la frase sobre la zona de prioridad. & SD8, 8.1.3; Tablas 8.2 a 8.4 \\
\LFXfiletefila
\mbox{SD8-12} & Dos «Tabla 8.4», Tabla 8.5 de siete columnas, fuente errada del RT-20 y RBS sin folio. & \textbf{Corregida.} Se unifica la numeración, se usan tablas de cinco columnas y se corrige la fuente RT-20. La RBS se actualiza a 36 riesgos, conservando folio y firma; se mantiene la separación del pie para evitar su superposición. & SD8, 8.1 a 8.3; Figura 8.1 (RBS) \\
\LFXfiletefila
\mbox{SD8-13} & Fuentes de las Tablas 8.2 a 8.4 y de la Declaración de IA con «contenido academico». & \textbf{Corregida.} Las fuentes pasan a «elaboración propia» y se elimina «contenido academico». La Declaración de IA explicita además el alcance de asistencia en los Anexos 8.C a 8.H; no constituye aceptación humana de esta actualización. & SD8, Tablas 8.2 a 8.4 y Declaración de IA, Anexos 8.C–8.H \\
\LFXfiletefila
\mbox{SD8-14} & Erratas: «tendria», «basandose», «academico», «asegurandose», «estmación». & \textbf{Corregida.} Se corrigen las erratas del cuerpo, T-16 y anexos, y se retira el punto final de los títulos de párrafo que producía puntuación duplicada. & SD8, 8.1.1; SD8-Anexos, justificación posterior a 8.F.2 y 8.H; T-16 \\
\LFXfiletefila
\mbox{SD8-15} & El riesgo del planificador no menciona la jubilación con fecha conocida ni la caída de 6 a 9 puntos. & \textbf{Corregida.} R-014 precisa jubilación dentro de dos años, sin fecha calendario especificada, y pérdida de 6 a 9 puntos porcentuales. Se solicita confirmar retiro y disponibilidad V-14 en mes 1, transferir en meses 1–3 y probar sin el planificador en meses 13–15. & T-16, R-014; SD8-Anexos, 8.A.2, 8.B y 8.D.1 \\
\LFXfiletefila
\mbox{SD8-16} & El SD13 asigna al piloto de R-034 meses que el Anexo 8.A no indica. & \textbf{Corregida.} CI-034 fija el paquete en los meses 11–15, con certificación del H5 en 11–12 y piloto real en 13–15. El SD13 remite al control CT-08 del Anexo 8.E con esas mismas ventanas. & SD8-Anexos, 8.E y 8.E.1; SD13, 13.1.4 y Anexo 13.B; T-19, ficha 1. \\
\LFXfiletefila
\mbox{SD8-17} & Falla mecánica, siniestro y asalto a camión: eventos operacionales del cliente de poco valor para el plan. & \textbf{Corregida.} Se mantienen R-021 a R-023 por la interrupción del reparto durante las pruebas. LafroX registra los casos TI afectados, preserva evidencia y repite las validaciones pendientes en un recorrido habilitado por el CLIENTE y el transportista. & SD8, 8.1; T-16, R-021 a R-023; SD8-Anexos, 8.A.1 y 8.F.1 \\
\LFXfiletefila
\mbox{SD8-18} & El T-16 menciona un «VME» que nunca se define ni se calcula. & \textbf{Corregida.} Se retira la mención. & T-16 \\
\LFXfiletefila
\mbox{SD8-19} & La Declaración de IA dice que la IA generó los riesgos y la tabla del T-16, con un solo revisor. & \textbf{Corregida.} La fila del T-16 describe asistencia de redacción y propuesta de valores. La declaración distribuye la revisión entre cinco responsables y amplía su alcance a 8.C–8.H. Se cotejó documentalmente esa declaración; falta confirmar la revisión humana de esta actualización, sin atribuirle aceptación técnica. & SD8, Declaración de IA, cuerpo, Anexos 8.C–8.H y T-16 \\
\LFXfiletefila
\mbox{SD8-20} & Estilo defensivo («no certifican», «no se presume», «no se afirma»). & \textbf{Corregida.} Se reemplazan salvedades repetidas por reglas de actuación y verificación. El cuerpo precisa quién coordina y qué evidencia verifica CAL. El T-16 explicita las condiciones para levantar bloqueos y cerrar SD8-06. & SD8; T-16 \\
\LFXfiletefila
\mbox{SD8-21} & Cuerpo de unas 14 páginas de contenido para un ítem del 10 \%. & \textbf{Corregida.} El contenido principal se amplía de 14 a 19 páginas con cinco cuadros de FMEA aplicado, respuestas prioritarias, costo-beneficio, reservas y condiciones de habilitación. Se explican las decisiones en el cuerpo y se conserva el detalle en anexos, sin alterar riesgos ni cálculos. & SD8, 8.2.3 y 8.3.1–8.3.3; Tablas 8.2.3-A, 8.3.1-A/B, 8.3.2-A y 8.3.3-A \\
\end{longtable}
}

\Needspace{8\baselineskip}
\subsection{Subdocumento 9: Plan de calidad (T-13 y T-17)}
{\linespread{1}\LFXcuerpotabla
\setlength{\parskip}{0pt}\renewcommand{\arraystretch}{1.10}\setlength{\tabcolsep}{4pt}
\begin{longtable}{F{0.09} F{0.29} F{0.40} F{0.22}}
\caption[Subdocumento 9: Plan de calidad (T-13 y T-17)]{Subdocumento 9: Plan de calidad (T-13 y T-17). Fuente: revisión del Informe 2 y versión corregida de LafroX.}\\[2pt]
\LFXtopruleoriginal \cab{N.º} & \cab{Observación} & \cab{Respuesta} & \cab{Sección modificada} \\ \midrule
\endfirsthead
\multicolumn{4}{@{}l@{}}{\fontsize{9pt}{11pt}\selectfont\color{lafroxGris}Tabla \thetable\ (continuación)}\\[3pt]
\LFXtopruleoriginal \cab{N.º} & \cab{Observación} & \cab{Respuesta} & \cab{Sección modificada} \\ \midrule
\endhead
\midrule
\multicolumn{4}{@{}r@{}}{\fontsize{9pt}{11pt}\selectfont\color{lafroxGris}Continúa en el folio siguiente.}\\
\endfoot
\bottomrule
\endlastfoot
\mbox{SD9-01} & “Falta una prueba explícita del manifiesto de identidad durante un corte (credenciales y relevos sin enlace), que el S4 declara como la solución de la identidad sin conexión; no aparece en el S9, en sus anexos ni en el T-13. AL-OFF-01 menciona «relevos» sin criterio de salida sobre la autenticación.” (p. 40) & \textbf{Corregida.} Se incorpora AL-ID-01 para H5/H10: emisión y relevo de credenciales sin enlace, expiración y revocación. La salida exige cero accesos indebidos y una cola íntegra y conciliada. Es un protocolo previsto, no una prueba ejecutada. & SD9, 9.2, Tabla 9.8; T-13, 7.5. Origen: SD4, 4.1 y Anexo 4-M. \\
\LFXfiletefila
\mbox{SD9-02} & “Usabilidad: se mide aprendizaje ≤ 2 h en bodega (págs. 10 y T-13, pág. 17), pero no se dice que la prueba se haga con personal recién ingresado al turno de noche, que es la condición de la restricción N° 11 (rotación del 38 \%). Debe declararse la población y el horario de la prueba.” “Declarar que la prueba de usabilidad de bodega se hace en turno de noche con personal de menos de un mes, y el umbral de productividad de la primera semana.” (pp. 40, 42) & \textbf{Corregida.} Se fija la prueba entre 22:00 y 06:00 con personal de bodega de menos de un mes. Criterios: aprendizaje ≤ 2 h, error ≤ 5 \% y al menos 80 \% del ritmo con papel en cada una de las primeras dos semanas. & SD9, 9.1; T-13, 7.3. \\
\LFXfiletefila
\mbox{SD9-03} & “Vocabulario inconsistente sobre el paso a producción: el S9 dice que el acta de cierre de la marcha blanca «habilita el paso a producción» y que G5 sólo abre el ambiente (págs. 16 y 30), mientras el T-13 habla de «el paso a producción de la Etapa 1, en el mes 13» (T-13, pág. 20). Debe usarse un solo término alineado con el T-18 de las Bases (producción desde los meses 16 y 21).” (p. 41) & \textbf{Corregida.} Un solo par de términos en SD9, T-13 y T-17: marcha blanca en el ambiente de Producción (M13–15 y M19–20, habilitada por G5 en H6 y H11) y paso a producción (M16 y M21, tras el acta de cierre, conforme al T-18). El acta integrada por etapa y la ventana de corte quedan fijadas en el acta de inicio 1.1.3. & SD9, 9.2–9.3; T-13, 7 y 9.2; T-17, 2.3, 4.1 y 4.3. Cotejo: T-15, 5.5; T-18. \\
\LFXfiletefila
\mbox{SD9-04} & “La acreditación CMMI se apoya en una evaluación Benchmark «n.º 58190, de noviembre de 2023» con renovación «antes de noviembre de 2026» (pág. 11), es decir, antes de la adjudicación; el propio texto advierte que no se invoca como garantía mientras no se acredite. Debe adjuntarse la evidencia o declararse como plan.” (pp. 38–39) & \textbf{Corregida.} La evaluación se conserva como antecedente histórico y la renovación se declara como plan pendiente. No se invoca como acreditación vigente; su estado se informa en H1. & SD9, 9.1. \\
\LFXfiletefila
\mbox{SD9-05} & “La restauración mensual de respaldos corre «todos los meses desde el 13» (pág. 27), incluidos septiembre y diciembre, sin decir sobre qué ambiente se restaura en esos meses.” (p. 40) & \textbf{Corregida.} En septiembre y diciembre la prueba de restauración usa Preproducción aislada y un respaldo sintético representativo. Valida el procedimiento y su tiempo, no la integridad de copias reales de Producción. & SD9, 9.3; T-13, 9.3. \\
\LFXfiletefila
\mbox{SD9-06} & “Errata de ensamblaje: «extbfCorrecciones en congelamiento» (S9, pág. 31), resto de un comando de formato que se perdió en la conversión. Una lectura de cierre la habría detectado.” (p. 41) & \textbf{Corregida.} El encabezado dice «Correcciones en congelamiento», sin el fragmento del comando LaTeX. & SD9, 9.3. \\
\LFXfiletefila
\mbox{SD9-07} & “Texto sobrecargado de salvedades defensivas que oscurecen la regla: «La segunda revisión nunca se presume instantánea» (pág. 30), «el máximo previsto no equivale a personas ya disponibles» (pág. 26), «No se garantiza que baste el paquete original ni se presupone una reserva de un riesgo concreto» (pág. 31).” (p. 41) & \textbf{Corregida.} Se sustituyen las salvedades por acciones: confirmar disponibilidad antes de reservar, escalar saldos negativos y tramitar capacidad adicional sin rebajar condiciones de aceptación. & SD9, 9.3. \\
\LFXfiletefila
\mbox{SD9-08} & Las pruebas ampliadas se concentran después de la construcción; el cronograma no refleja con suficiente detalle las pruebas durante la construcción (pp. 40–41). & \textbf{Corregida.} Durante la construcción y en cada despliegue en QA se ejecutan pruebas unitarias y de contrato. Las pruebas ampliadas corresponden a la certificación. & SD9, 9.3 y Anexo 9.E.1; T-13; T-15, 6 \\
\LFXfiletefila
\mbox{SD9-09} & “El T-13 dedica unas cincuenta páginas (sección 8.2, págs. 21–70) a la matriz caso por caso. Es correcto como formulario, pero conviene un resumen por módulo al inicio de la sección.” (p. 41) & \textbf{Corregida.} Se agrega un resumen por módulo al inicio, sin quitar la matriz detallada. Distingue los 1.328 casos previstos de cualquier resultado de ejecución. & T-13, 8.2. \\
\LFXfiletefila
\mbox{SD9-10} & “Tensión entre dos párrafos sobre cambios en congelamiento: la pág. 29 describe un procedimiento de cambio de emergencia en septiembre con autorización de la Contraparte Técnica; la pág. 31 dice que «Las Bases no prevén excepciones… y este plan no las presume». Debe quedar una sola regla.” (p. 41) & \textbf{Corregida.} Regla única: del 1 al 25 de septiembre se contiene y corrige en Preproducción. Desde el 26, solo en marcha blanca E2 se despliega una corrección de emergencia con autorización escrita y fuera del despacho. No hay paso a producción con incidentes altos o críticos abiertos. & SD9, 9.3.2; T-13, 9.2. Origen: T-18. \\
\LFXfiletefila
\mbox{SD9-11} & “En la Figura 9.4 el rótulo «Recuperación» se superpone a su recuadro (pág. 18).” (p. 41) & \textbf{Corregida.} Se separa «Recuperación» del borde del recuadro; se verificó visualmente en el PDF. & SD9, Figura 9.4. \\
\LFXfiletefila
\mbox{SD9-12} & “Tres hitos quedan con holgura negativa una vez descontados los diez días de subsanación: H1 (−5), H7 (−8) y H12 (−7) (Tabla 9.13, pág. 30). El S9 lo reconoce y propone revisar los expedientes con la Contraparte Técnica antes de la entrega formal, pero no agrega tiempo en la Gantt: es una mitigación declarada, no un ajuste del plan.” (p. 41) & \textbf{Corregida.} H1, H7 y H12 se entregan por incrementos: 4 de marzo de 2027, y 5 de abril y 5 de septiembre de 2028. El incremento final se subsana en cinco días hábiles, con el Líder de Calidad y un Analista QA Senior asignados. La holgura tras la subsanación pasa a 0, +5 y +3 días hábiles, sin mover las fechas nominales ni los meses del T-15. & SD9, 9.3, Tabla 9.13; T-17, 2.3, 2.4 y 4.1. Cotejo: T-15, 5.5. \\
\LFXfiletefila
\mbox{SD9-13} & “Parcial: matriz de trazabilidad. El T-12 lleva cada requerimiento a componente, paquete EDT y sección, pero la prueba queda en el T-13 y el criterio en el T-17 (Anexos, folio 176), de modo que la cadena completa que exige el Cap. 17.1 (origen, requerimiento, componente, EDT, prueba, criterio) no se lee en una sola matriz.” (p. 10) & \textbf{Corregida.} El T-13 conecta origen, requisito, componente, EDT, prueba prevista y criterio con su fuente en una sola fila. Se verificaron 644 requisitos, 1.328 casos previstos y 22 brechas explícitas; no se presentan como pruebas ejecutadas. & T-13, 8.2. Cruce: SD3/T-12 (observación de Alcance) y T-17. \\
\end{longtable}
}

\Needspace{8\baselineskip}
\subsection{Subdocumento 13: Innovaciones (T-19)}
{\linespread{1}\LFXcuerpotabla
\setlength{\parskip}{0pt}\renewcommand{\arraystretch}{1.10}\setlength{\tabcolsep}{4pt}
\begin{longtable}{F{0.09} F{0.29} F{0.40} F{0.22}}
\caption[Subdocumento 13: Innovaciones (T-19)]{Subdocumento 13: Innovaciones (T-19). Fuente: revisión del Informe 2 y versión corregida de LafroX.}\\[2pt]
\LFXtopruleoriginal \cab{N.º} & \cab{Observación} & \cab{Respuesta} & \cab{Sección modificada} \\ \midrule
\endfirsthead
\multicolumn{4}{@{}l@{}}{\fontsize{9pt}{11pt}\selectfont\color{lafroxGris}Tabla \thetable\ (continuación)}\\[3pt]
\LFXtopruleoriginal \cab{N.º} & \cab{Observación} & \cab{Respuesta} & \cab{Sección modificada} \\ \midrule
\endhead
\midrule
\multicolumn{4}{@{}r@{}}{\fontsize{9pt}{11pt}\selectfont\color{lafroxGris}Continúa en el folio siguiente.}\\
\endfoot
\bottomrule
\endlastfoot
\mbox{SD13-01} & “Dimensionar el beneficio de cada innovación en unidades físicas (unidades vencidas, lotes, horas, entregas) a partir de los datos del caso, sin montos.” (p. 44) · INN-01 & \textbf{Corregida.} Escenario exploratorio: 6.000 unidades vencidas evitadas al mes desde el mes 16 (2,4 millones de unidades × 1 \% × 25 \%, con tasa y reducción supuestas). Se contrasta en el mes 28 por ruta, local, lote y período. No se suma a merma ni a INN-03. & SD13, 13.1.5 y Anexo 13.C; T-19, ficha 1, campo 12. Cotejo: modelo financiero, Innovaciones. \\
\LFXfiletefila
\mbox{SD13-02} & “Innovación 3 (tipo 3): técnicamente la más sólida y la de mejor sustento bibliográfico. Pero el T-12 marca el RT-05.30 (deseable) como «Cumple» mediante esta misma innovación y el S5 lo «oferta» con ella (pág. 27; S5, pág. 34): la misma capacidad se cuenta como cumplimiento de un requisito de las Bases y como innovación.” (pp. 43–44) & \textbf{Corregida.} El RT-05.30 queda cubierto por un modelo cinético determinista. La innovación es el uso operativo adicional: cálculo por lote sin enlace, reconstrucción central y uso del resultado por Calidad. & SD13, 13.3.2; T-19, ficha 3. Arrastre: SD3/T-12 y SD5, 5.2.6 y Anexo A.20. \\
\LFXfiletefila
\mbox{SD13-03} & “Innovación 2 (tipo 2): bien argumentada frente a los reportes de fallos y al RT-14.06 (pág. 18), pero el proceso que cambia es el de diagnóstico del propio adjudicatario y su beneficio son «horas de diagnóstico y regresiones repetidas evitadas» (pág. 23). Para Puelche el efecto es indirecto, y la reproducción de incidentes con entornos efímeros es una práctica de ingeniería del proveedor. Debe defender en la presentación por qué es una innovación de proceso del CLIENTE.” (p. 43) & \textbf{Corregida.} La persona de terreno identifica la tarea interrumpida y confirma que se reanudó; Puelche registra aviso y reanudación. El proveedor sigue a cargo de reproducir y corregir el incidente. El beneficio físico se sigue por separado (SD13-07). & SD13, 13.2 y Anexo 13.C; T-19, ficha 2. \\
\LFXfiletefila
\mbox{SD13-04} & “La introducción dice que la valorización «se desarrolla para el Informe 3» (pág. 7): en una oferta se remite a la Oferta Económica, no a una instancia del proceso.” (p. 44) & \textbf{Corregida.} La valorización se remite a la Oferta Económica, Entregable 2, ítem 2.7; la oferta técnica no incorpora los montos. & SD13, introducción y apartados 13.1.5 a 13.5.5; Anexo 13.B. \\
\LFXfiletefila
\mbox{SD13-05} & “La Declaración de uso de IA indica que la revisión humana «se centró en la redacción, claridad y comprensión del contenido» (pág. 50), con un solo revisor para todo el capítulo y nivel «Alto» en todas las figuras. No acredita revisión técnica de los modelos, umbrales ni cálculos; en la defensa deben explicarlo quienes diseñaron cada innovación.” (p. 44) & \textbf{Pendiente.} La declaración distingue la revisión editorial de la técnica. El Anexo 13.C define seis fichas (cinco innovaciones y la Figura 13.1), con revisor, fecha, versión revisada, elemento verificado, evidencia y resultado. La revisión humana verificará las correcciones aplicables. Se corrige para la Propuesta Final. & SD13, Declaración de uso de IA y Anexo 13.C, pauta de revisión; T-19. \\
\LFXfiletefila
\mbox{SD13-06} & “Las fechas de los paquetes remiten al S7, T-14, T-15, T-16 y Anexo 8.A, que no forman parte de este cruce; deben verificarse con los capítulos nuevos.” (p. 44) & \textbf{Corregida.} H4 se entrega el 20-10-2027 con los módulos probados, y el 22 se agrega el último incremento. H5 y H10 se entregan completos el 24-11-2027 y el 27-06-2028, con las pruebas terminadas el 19-11 y el 15-06. El paquete cierra al preparar el acta (1-12 y 28-06). Las fechas y los meses de cobro del T-15 no cambian. & SD13, Anexo 13.B.1; T-19; T-17, 2.4. Cotejo: T-14/T-15 y SD8, Anexo 8.E/T-16. \\
\LFXfiletefila
\mbox{SD13-07} & “Dimensionar el beneficio de cada innovación en unidades físicas (unidades vencidas, lotes, horas, entregas) a partir de los datos del caso, sin montos.” (p. 44) · INN-02 & \textbf{Corregida.} Escenario exploratorio: 156 minutos de interrupción evitados al mes (312 minutos-persona) desde el mes 13, con un incidente por semana, dos personas, 120 minutos y −30 \% como supuestos. Se contrasta desde el mes 15, del aviso a la reanudación confirmada; I-02A mide diagnóstico. & SD13, 13.2.5 y Anexo 13.C; T-19, ficha 2, campo 12. Cotejo: modelo financiero, Innovaciones. \\
\LFXfiletefila
\mbox{SD13-08} & “Dimensionar el beneficio de cada innovación en unidades físicas (unidades vencidas, lotes, horas, entregas) a partir de los datos del caso, sin montos.” (p. 44) · INN-03 & \textbf{Corregida.} Escenario exploratorio: 1,5 lotes menos afectados por vencimiento al mes desde el mes 16 (150,6 lotes de familias habilitadas al mes × 5 \% × 20 \%, o 10 por cada 1.000). En el mes 28, el conteo de lotes únicos por familia reemplaza a la ponderación. No se suma a INN-01 ni a merma. & SD13, 13.3.5 y Anexo 13.C; T-19, ficha 3, campo 12. Cotejo: modelo financiero, Innovaciones. \\
\LFXfiletefila
\mbox{SD13-09} & “Dimensionar el beneficio de cada innovación en unidades físicas (unidades vencidas, lotes, horas, entregas) a partir de los datos del caso, sin montos.” (p. 44) · INN-04 & \textbf{Corregida.} Escenario exploratorio: 31.000 minutos-persona menos al mes desde el mes 24 (31.000 pedidos × 20 minutos × 5 \%), es decir, 16 h 40 min por cada 1.000 entregas. No son HH presupuestadas; se contrasta por actividad y régimen sin degradar el servicio ni el OTIF. & SD13, 13.4.5 y Anexo 13.C; T-19, ficha 4, campo 12. Cotejo: modelo financiero, Innovaciones. \\
\LFXfiletefila
\mbox{SD13-10} & “Dimensionar el beneficio de cada innovación en unidades físicas (unidades vencidas, lotes, horas, entregas) a partir de los datos del caso, sin montos.” (p. 44) · INN-05 & \textbf{Corregida.} Escenario exploratorio: 101.296 minutos del almacenero evitados al mes desde el mes 21 (50.648 decisiones × 2 minutos; 62.000 visitas ponderadas por el canal tradicional). El efecto comercial se valoriza solo si se demuestra al mes 33. & SD13, 13.5.5 y Anexo 13.C; T-19, ficha 5, campo 12. Cotejo: modelo financiero, Innovaciones. \\
\LFXfiletefila
\mbox{SD13-11} & “Incorporar las cinco innovaciones al alcance del Subdocumento 3 (3.1 y 3.4)” (p. 45) & \textbf{Corregida.} Alcance incorpora las cinco innovaciones y remite al SD13 y al T-19; se conserva la distinción entre el requisito analítico y el aporte innovador. Véase SD3-14. & SD3, 3.1 y 3.4; remisión a SD13 y T-19. \\
\LFXfiletefila
\mbox{SD13-12} & “Completar la red: toda la migración, las integraciones con las cadenas y las fases de cada innovación deben tener predecesoras; recalcular el CPM a partir de duraciones y dependencias, no de fechas fijadas.” (pp. 29, 32) & \textbf{Corregida.} La red del T-15 vincula las fases de las cinco innovaciones mediante siete enlaces fin–inicio y entregas habilitantes; 7.4.1 y 7.4.2 preceden a 7.4.3. El CPM se calcula con duraciones y dependencias. El T-15 programa 7.1.4, 7.2.4 y 7.3.4 hasta el 28 de abril de 2028, fecha de fin tardío del CPM. Las sesiones EV-15-01 a EV-15-03, del 3 al 28 de abril, tienen 48 HH cada una, con desglose por perfil y evidencia de cierre prevista. Las 144 HH forman parte de las 533 HH de Datos y las 154 HH de Calidad programadas para abril. & SD13, Anexo 13.B y 13.B.1; vista editable de Innovaciones; T-19, fichas 1–3. Cotejo: T-15; modelo financiero, HH por rol y mes. \\
\end{longtable}
}
